% !TEX root = ../main.tex
% ---------------------------------------------------------------
% GENERATED FILE - DO NOT EDIT.
% Produced by docs/tools/render_reference.py from the repository source.
% Regenerate with: make docs
% ---------------------------------------------------------------

\apibreak
\section{\texttt{GuardrailDecision}}
\label{class:feedback_intelligence_agent.guardrails.GuardrailDecision}\label{schema:GuardrailDecision}\index{Classes!GuardrailDecision}

Outcome of one guardrail evaluation.
\apifield{Inherits}\texttt{BaseModel}
\apifield{Fields}
\begin{arglist}
  \item[\texttt{allowed}] \texttt{bool}, required. Whether the request may proceed.
  \item[\texttt{reason}] \texttt{str}, required. Human-readable explanation, including the rule that fired.
  \item[\texttt{severity}] \texttt{Severity}, default \texttt{'low'}. \texttt{low} (benign or malformed), \texttt{medium} (policy bypass attempts), or \texttt{high} (injection or data exfiltration attempts).
  \item[\texttt{suggested\_\allowbreak{}response}] \texttt{str \textbar{} None}, default \texttt{None}. Safe text the caller can return verbatim when the request is blocked.
\end{arglist}
\apifield{Source}\texttt{src/\allowbreak{}feedback\_\allowbreak{}intelligence\_\allowbreak{}agent/\allowbreak{}guardrails.\allowbreak{}py:36}~(\href{https://github.com/DiogoRibeiro7/feedback-intelligence-agent/blob/b6cbc710ffdfa5f26f700c999d594c67588c1e16/src/feedback_intelligence_agent/guardrails.py#L36}{online}), pydantic model, module \cref{module:feedback_intelligence_agent.guardrails}
